\documentclass[11pt]{article}
\usepackage[margin=1in]{geometry}
\usepackage[T1]{fontenc}
\usepackage[utf8]{inputenc}
\usepackage{lmodern}
\usepackage{amsmath,amssymb,mathtools,booktabs,graphicx,natbib,hyperref}

\title{Decision-Aware Penalized Trend Estimation for Financial Time Series}
\author{Heriberto Espino Montelongo}
\date{\today}

\newcommand{\R}{\mathbb{R}}
\newcommand{\Dt}{D_d}

\begin{document}
\maketitle

\begin{abstract}
We study the selection of smoothness in finite-difference penalized trend
estimators when the fitted trend is used for temporal forecasting and, in a
later stage, downstream financial decisions. The software framework separates
a pure quadratic smoother from a Guerrero-style formulation with an estimated
difference drift, because their dependence on the penalty parameter is not the
same. The pure model admits explicit first- and second-order sensitivities with
respect to the penalty parameter, enabling numerical optimization in
log-penalty space. The empirical design is strictly chronological and compares
statistical selection criteria with forecast-based and, eventually,
decision-based objectives. This manuscript is a research draft: no novelty or
performance claims are made until the corresponding comparisons are completed.
\end{abstract}

\section{Problem formulation}
Let $y\in\R^n$ and let $\Dt$ be the order-$d$ finite-difference matrix. Define
$Q=\Dt^\top\Dt$.

\subsection{Pure penalized trend}
The first model is
\begin{equation}
\widehat t_{\lambda,d}
=\arg\min_{t\in\R^n}\left\{\|y-t\|_2^2+\lambda\|\Dt t\|_2^2\right\},
\qquad \lambda\ge 0,
\end{equation}
with solution
\begin{equation}
\widehat t_{\lambda,d}=(I+\lambda Q)^{-1}y.
\end{equation}
Writing $S_\lambda=(I+\lambda Q)^{-1}$ gives
\begin{align}
\frac{\partial \widehat t}{\partial\lambda}
&=-S_\lambda Q\widehat t,\\
\frac{\partial^2 \widehat t}{\partial\lambda^2}
&=2S_\lambda Q S_\lambda Q\widehat t.
\end{align}
These derivatives are evaluated numerically without constructing an explicit
matrix inverse.

\subsection{Guerrero-style trend with difference drift}
The second model is
\begin{equation}
\widehat\tau_{\lambda,d}
=(I+\lambda Q)^{-1}
\left(y+\lambda\widehat m(\lambda)\Dt^\top\mathbf 1\right).
\end{equation}
Because $\widehat m$ is re-estimated from the fitted trend,
$\widehat m=\widehat m(\lambda)$. Consequently the pure-model derivative cannot
be reused unchanged. A full sensitivity calculation must differentiate the
coupled fixed-point/system; implicit differentiation is the intended route.

\section{Chronological model selection}
Every validation protocol respects time ordering. At forecast origin $T$, only
$y_1,\ldots,y_T$ may enter the fit. Future observations may score the forecast
but may not influence the trend being evaluated. We will compare grid search,
Brent-type minimization in $\theta=\log\lambda$, Newton iterations using
analytic derivatives where available, generalized cross-validation, and
rolling-origin forecast loss.

\section{Research hypothesis}
The central empirical question is whether a penalty selected for statistical
or forecasting accuracy is also optimal for a downstream financial objective.
Symbolically,
\begin{equation}
\lambda_{\mathrm{forecast}}\overset{?}{=}\lambda_{\mathrm{decision}}.
\end{equation}
This equality is treated as a testable hypothesis, not an assumption.

\section{Planned experiments}
Experiments will first validate derivatives on synthetic data against centered
finite differences. They will then compare the two smoothers and classical
baselines under rolling-origin evaluation. Only after those components are
stable will the trend signal be connected to a portfolio layer.

\section{Reproducibility}
All reusable estimators live in \texttt{src/trend\_estimation}; manuscript
experiments must call the installed package rather than duplicate estimator
logic. The companion tutorial manuscript documents each derivation and
implementation decision at a slower pedagogical level.

\bibliographystyle{plainnat}
\IfFileExists{paper/biblio.bib}{\bibliography{paper/biblio}}{\bibliography{../biblio}}
\end{document}
